% Standalone schematic — inverting op-amp (classic negative feedback)
% Compile: xelatex negative_feedback.tex   (or pdflatex)
% Embed:  \includegraphics{schmatic/negative_feedback}
\documentclass[border=8pt]{standalone}
\usepackage{circuitikz}
\usetikzlibrary{calc}
\ctikzset{bipoles/length=1.3cm}

\begin{document}
\begin{circuitikz}[american, font=\small, line width=0.7pt]

  % ---- Op-amp ----
  \node [op amp, yscale=1] (OA) at (0,0) {};

  % ---- Inverting input summing junction ----
  \coordinate (A) at ($(OA.-) + (-0.6,0)$);
  \draw (OA.-) -- (A);
  \node[circ] at (A) {};

  % ---- Input resistor R1 ----
  \draw (A) to[R, l_=$R_1$, -o] ++(-3.0,0)
        node[left] {$v_\mathrm{in}$};

  % ---- Feedback resistor R_f (top path) ----
  \coordinate (B) at ($(A)        + (0,1.7)$);
  \coordinate (C) at ($(OA.out)   + (0.7,1.7)$);
  \draw (A) -- (B)
        to[R, l=$R_f$] (C);

  % ---- Output node ----
  \coordinate (D) at ($(OA.out) + (0.7,0)$);
  \node[circ] at (D) {};
  \draw (OA.out) -- (D) -- (C);
  \draw (D) to[short, -o] ++(1.4,0)
        node[right] {$v_\mathrm{out}$};

  % ---- Non-inverting input to ground ----
  \draw (OA.+) -- ++(-0.6,0) node[ground] {};

\end{circuitikz}
\end{document}
